% &%%%%%%%%%%%%%%%%%%%%%%%%%%%%%%%%%%%%%%%%%%%%%%%%%%%%%%%%%%%
% &%%                                                       %%
% &%%                     INTRODUCTION                      %%
% &%%                                                       %%
% &%%%%%%%%%%%%%%%%%%%%%%%%%%%%%%%%%%%%%%%%%%%%%%%%%%%%%%%%%%%

% !===========================================================

\graphicspath{{5_Images/0_Introduction/}}                       % Chemin vers le stockage des images du chapitre
\makeatletter
\def\input@path{{5_Images/0_Introduction/}}                     % Pareil que au dessus
\makeatother

% !===========================================================

\chapter*{Introduction}\label{Chap:Introduction}                % Ajout du chapitre en tant que chapitre fantôme (pour ne pas lui attribuer de numéro)

\mtcaddchapter%                                                 % Permet d'ajouter manuellement un mini-ToC fantôme à cause du \chapter*
\addcontentsline{toc}{chapter}{Introduction}                    % Ajout du chapitre à la ToC
\markboth{Introduction}{Introduction}                           % Permet d'afficher le chapitre en en-tête

% !===========================================================

% Je vous montre un exemple de la structure de mon introduction mais adaptez à votre sauce

% ?-----------------------------------------------------------
\section*{Préambule}
% ?-----------------------------------------------------------

\lettrine{L}{orem} ipsum...

\lipsum[1-3]%

% ?-----------------------------------------------------------
\section*{Contexte industriel}
% ?-----------------------------------------------------------

\lipsum[1-3]%

% ?-----------------------------------------------------------
\section*{Cadre technique et scientifique}
% ?-----------------------------------------------------------

\lipsum[1-3]%